\section{Related Work}
\label{sec:related}

\textbf{mmWave point-cloud HAR.} RadHAR \cite{singh2019radhar} released the
MMActivity data used here. It voxelised each frame into a $10\times32\times32$
grid and classified sequences with a time-distributed CNN and a BiLSTM (90.47\%
accuracy). Later work moved to models that consume points directly.
MMPoint-GNN \cite{gong2021mmpointgnn} builds dynamic graphs per frame. Lee and
Kim \cite{lee2022gnn} pre-train a GNN on 3-D joint coordinates and report
96.55\% on RadHAR. Gu et al.\ \cite{gu2024lpn} pair a light PointNet with a
bidirectional linear-attention LSTM (95.12\%, 79.7\,K parameters).
MAK-GCN \cite{zakka2025makgcn} learns multi-head adaptive graph kernels (97.54\%,
1.86\,M parameters). OG-PCL \cite{yan2025ogpcl} uses occupancy-gated
convolutions (91.75\%, 0.83\,M parameters) and reports PCT+LSTM
\cite{guo2021pct} at 92.78\%. Gao et al.\ \cite{gao2025star} model each frame as
a star graph around a virtual centre. Lim et al.\ \cite{lim2024fpga} deploy a
4-bit pillar-based classifier on an FPGA. For gestures, Pantomime
\cite{palipana2021pantomime} and Tesla-Rapture \cite{salami2023tesla} showed
that sparse radar point clouds support lightweight recognisers. The same sensing
also drives rehabilitation \cite{an2021mars}, mesh reconstruction
\cite{xue2021mmmesh,chen2022mmbody} and scene flow \cite{ding2024milliflow}.
Ullmann et al.\ \cite{ullmann2023survey} survey radar-based continuous HAR.

\textbf{Benchmarks and evaluation protocols.} MiliPoint \cite{cui2023milipoint}
(49 actions, 11 subjects) and MM-Fi \cite{yang2023mmfi} are recent large-scale
benchmarks. On MiliPoint action classification the original authors report at most
34.45\% (PointNet++), whereas MAK-GCN reports 98.25\% under an
80/10/10 split. Several RadHAR studies validate on 20\% of training windows drawn at
random, and those windows overlap by 50 of 60 frames. Data leakage of this kind is a
known way to inflate results \cite{kaufman2012leakage}. To our knowledge, no
prior work has measured its effect on radar point-cloud HAR.

\textbf{Point-cloud backbones.} PointNet \cite{qi2017pointnet} and PointNet++
\cite{qi2017pointnet2} introduced learning on raw point sets. DGCNN
\cite{wang2019dgcnn} builds dynamic $k$-NN graphs. Point Transformer
\cite{zhao2021pt}, PCT \cite{guo2021pct}, PTv2/PTv3 \cite{wu2022ptv2,wu2024ptv3},
PointNeXt \cite{qian2022pointnext} and PointMLP \cite{ma2022pointmlp} raised accuracy
on dense LiDAR and CAD data. P4Transformer \cite{fan2021p4t} handles point-cloud
videos. All of these assume hundreds to thousands of points per frame and rely on
FPS and $k$-NN, which are costly on micro-controllers.

\textbf{Research gap.} (i) Most radar HAR models either reuse dense-point
backbones or voxelise the data. Few exploit what is specific to radar: extreme
sparsity, a direct Doppler measurement, and static clutter. (ii) Reported
accuracies are rarely obtained with leakage-free (recording- or subject-disjoint)
splits, and protocols differ between papers. (iii) Cross-dataset transfer
between radar datasets and edge cost (latency, size, quantisation) are seldom
reported together. This work addresses all three gaps.
